% ============================================================
% V. DISCUSIÓN
% ============================================================

\section{Discusión}

Los resultados muestran que la Transformada Rápida de Fourier constituye una herramienta útil para analizar señales de audio y observar diferencias entre diferentes fuentes sonoras. La frecuencia dominante fue, de las tres características calculadas, la que presentó mayor utilidad para distinguir las muestras analizadas de forma cualitativa: en el caso de los animales, las frecuencias dominantes estuvieron suficientemente separadas en términos absolutos para observar diferencias claras entre las tres señales.

Sin embargo, el puntaje de separabilidad de la Sección~\ref{sec:resultados-separabilidad} matiza esa lectura: al normalizar cada diferencia de frecuencia dominante por la dispersión de amplitud combinada de las dos señales comparadas ($\sigma_a + \sigma_b$), todos los pares evaluados quedan por debajo del umbral de separabilidad parcial. Es decir, una diferencia de frecuencia que en Hz absolutos parece grande (por ejemplo, los $1449.90$ Hz entre gato y perro) sigue siendo pequeña frente a la variabilidad de amplitud reportada para esas mismas señales. Esto ilustra por qué comparar únicamente frecuencias dominantes, sin considerar su dispersión, puede sobreestimar cuán separables son en la práctica dos fuentes sonoras.

En el caso de los instrumentos, el contrabajo presentó una frecuencia dominante considerablemente menor que el piano y la flauta, pero la diferencia entre piano y flauta fue menor, lo que demuestra que la clasificación de un instrumento como grave, medio o agudo depende de la nota o muestra específica utilizada; esto se refleja en que el par Piano--Flauta obtuvo el puntaje de separabilidad más bajo de toda la Tabla~\ref{tab:separabilidad} ($S = 0.01$).

Por otra parte, la media de las señales presentó valores cercanos a cero en la mayoría de los casos (con la excepción de Persona 1, cuyo valor de $6.699$ resulta atípico frente al resto de la tabla y debería verificarse, por ejemplo revisando si esa grabación tiene un componente de continua o \emph{offset} de captura). Esto indica que la media, tal como se calcula aquí, no resulta por sí sola adecuada para diferenciar las fuentes sonoras.

La desviación estándar presentó diferencias entre las señales, pero no mostró una relación directa con la frecuencia dominante: el contrabajo y la flauta presentaron desviaciones estándar similares a pesar de tener características espectrales distintas. Cabe señalar que \texttt{animal\_analyzer.core.signal\_stats} calcula esta desviación estándar sobre las muestras devueltas por \texttt{librosa.load}, las cuales están típicamente normalizadas al rango $[-1, 1]$; para una señal de audio real en ese rango se esperarían valores de $\sigma$ del orden de $10^{-1}$, no del orden de $10^{3}$ como los reportados en la Tabla~\ref{tab:resultados}. Esta discrepancia de escala sugiere que los valores de esta tabla provienen de una etapa de procesamiento previa a la normalización de amplitud (por ejemplo, muestras PCM enteras sin escalar) y no afecta a la frecuencia dominante (invariante ante un reescalado de amplitud), pero sí infla artificialmente el denominador de la Ec.~\ref{eq:separabilidad} y, por tanto, reduce los puntajes de separabilidad de la Tabla~\ref{tab:separabilidad}. Se recomienda re-ejecutar el análisis pasando estas mismas grabaciones por el pipeline unificado (\texttt{load\_audio} + \texttt{analyze\_audio}) antes de tomar los puntajes de separabilidad como definitivos.

Debe considerarse además que la magnitud de la FFT depende de factores como la amplitud de grabación, la duración de la señal y las condiciones de captura, por lo que no debe utilizarse directamente como criterio de identificación sin una normalización previa consistente entre grabaciones.

Finalmente, aunque se observaron diferencias entre las señales, el número reducido de muestras (una grabación por categoría, dos para voces) impide determinar si dichas diferencias son estadísticamente significativas o generalizables. Para desarrollar un sistema de identificación más robusto sería necesario ampliar \texttt{audio\_samples/} con múltiples grabaciones por categoría y considerar características adicionales a la frecuencia dominante y la dispersión global, como la distribución de armónicos o la energía por bandas.
